\documentclass{amsart}
\usepackage{amsmath,amssymb,amsthm,hyperref}
\usepackage{algorithm}
\usepackage{algpseudocode}
\newtheorem{theorem}{Theorem}
\newtheorem{lemma}{Lemma}
\newtheorem{proposition}{Proposition}
\newtheorem{corollary}{Corollary}
\theoremstyle{definition}
\newtheorem{definition}{Definition}
\newtheorem{remark}{Remark}

\begin{document}

\title{Geodetic homeomorphs of the Hoffman--Singleton graph: a complete
characterization and enumeration}
\maketitle

\section{Statement and conventions}

Let $H$ be the Hoffman--Singleton graph: the unique $(7,5)$--Moore graph, which
is $7$--regular on $50$ vertices, has $|E(H)|=175$ edges, girth $5$, diameter
$2$, and is geodetic (between any two vertices there is a unique shortest path).
For a Moore graph of diameter $2$ this is automatic: adjacent vertices have no
common neighbour (no triangles), and non-adjacent vertices have exactly one
common neighbour, so every shortest path is unique.

For a positive-integer assignment $x=(x_e)_{e\in E(H)}\in\mathbb Z_{>0}^{175}$
let $G'=G'(x)$ be the subdivision of $H$ in which each edge $e=uv$ is replaced
by an internally disjoint path of length $x_e$ joining $u$ and $v$. The $50$
original vertices, of degree $7$, are the \emph{branch vertices}; the new
degree-$2$ vertices are \emph{internal}. We call $G'(x)$ a \emph{homeomorph} of
$H$. We must decide for which $x$ the graph $G'(x)$ is geodetic, enumerate all
such $x$, and give an algorithm doing so.

Throughout, an \emph{$H$--path} between branch vertices $u,v$ is a path in $H$
from $u$ to $v$; its \emph{weight} is $\sum_{e}x_e$ over its edges. A
\emph{pentagon} (resp.\ \emph{hexagon}) is a $5$--cycle (resp.\ $6$--cycle) of
$H$; $H$ has exactly $1260$ pentagons and $5250$ hexagons. For a cycle $C$ and
two of its vertices $a,b$, the two $a$--$b$ paths along $C$ are the \emph{arcs}
of $C$ between $a$ and $b$, with weights equal to the sums of $x_e$ along them.
If the shorter arc (in number of edges) uses $k$ edges and the longer uses
$\ell-k$ (where $\ell$ is the length of the cycle), we call this the
\emph{$k{+}(\ell-k)$ split}.

We first record the structural facts used repeatedly. They are classical
(Hoffman--Singleton 1960) and are verified directly from the Robertson
pentagon/pentagram model in the accompanying computations (\texttt{check.py};
$50$ vertices, $175$ edges, $7$-regular, triangle-free, diameter $2$,
$1260$ pentagons, $5250$ hexagons, all confirmed).

\begin{lemma}[Structure of $H$]\label{lem:struct}
$H$ is $7$--regular on $50$ vertices, $|E(H)|=175$, has girth $5$ and diameter
$2$, is triangle-free, and is geodetic. Consequently:
\begin{enumerate}
\item every pair of adjacent vertices lies on no triangle, so the shortest
$H$--path between them other than the edge has at least $4$ edges;
\item every pair of non-adjacent vertices $u,v$ has a unique common neighbour
$w$, and the unique shortest $H$--path between them is $u\,w\,v$.
\end{enumerate}
\end{lemma}

\section{Distances in $G'(x)$ reduce to weighted $H$--paths}

\begin{lemma}\label{lem:distbranch}
For branch vertices $u,v$,
$d_{G'(x)}(u,v)=\min_{P}\ \sum_{e\in P}x_e$, the minimum over all $H$--paths $P$
from $u$ to $v$. Moreover every $u$--$v$ path in $G'(x)$ either lies inside a
single subdivided edge (only when $u,v$ are adjacent in $H$) or passes through a
sequence of branch vertices $u=b_0,b_1,\dots,b_k=v$ with $b_ib_{i+1}\in E(H)$,
hence projects to an $H$--walk; a shortest one projects to an $H$--path. Distinct
$u$--$v$ paths of $G'(x)$ that meet only at $u,v$ project to distinct
$H$--paths.
\end{lemma}

\begin{proof}
Every internal vertex of $G'(x)$ has degree $2$; thus any walk in $G'(x)$, on
entering a subdivided edge $e=ab$ at one endpoint, must traverse it entirely to
the other endpoint before leaving (it cannot branch). Hence every $u$--$v$ walk
between branch vertices is determined by the sequence of branch vertices it
visits, an $H$--walk, with $G'$--length $\sum x_e$ over traversed edges. A
minimum-length walk repeats no branch vertex (deleting a closed sub-walk
shortens it), so it is an $H$--path of weight $\sum_{e}x_e$. Taking the minimum
gives the formula.
\end{proof}

We also need distances to internal vertices. Let $m$ be the internal vertex on
the subdivided edge $e=ab$ at $G'$--distance $s$ from $a$ ($0<s<x_e$). By the
degree-$2$ rigidity inside $e$ together with Lemma~\ref{lem:distbranch},
\begin{equation}\label{eq:internaldist}
d_{G'}(u,m)=\min\big(d_{G'}(u,a)+s,\ d_{G'}(u,b)+(x_e-s)\big)
\end{equation}
for any branch vertex $u$ (and analogously, with four combinations, when the
source is itself an internal vertex), because any path to $m$ must enter the
edge $e$ from $a$ or from $b$.

\section{The characterization theorem}

\begin{definition}\label{def:conditions}
For $x\in\mathbb Z_{>0}^{175}$ define:
\begin{itemize}
\item[\textnormal{(P)}] (\emph{pentagon parity}) every pentagon $C$ of $H$ has
\emph{odd} total weight $\sum_{e\in C}x_e$;
\item[\textnormal{(H)}] (\emph{geodesic hexagon non-tie}) for every hexagon $C$
of $H$ and every two of its branch vertices $a,b$ whose split is $1{+}5$ or
$2{+}4$, the two arcs of $C$ between $a,b$ have \emph{distinct} weights.
\end{itemize}
\end{definition}

\begin{remark}[Why only the $1{+}5$ and $2{+}4$ splits, not $3{+}3$]
\label{rem:why33excluded}
The $3{+}3$ split is deliberately \emph{excluded} from (H). By
Lemma~\ref{lem:cycles}(2) below, the two endpoints $a,b$ of a $3{+}3$ split are
at $H$--distance $2$ with their unique common neighbour lying \emph{off} the
hexagon, so neither $3$--arc is an $H$--geodesic. Equal weights of the two
$3$--arcs therefore produce two equal-weight \emph{non}-shortest paths, which
does not violate geodeticity. Indeed the all-ones assignment ($G'=H$, certainly
geodetic) and every uniform odd assignment $x_e\equiv c$ have both $3$--arcs of
each hexagon equal to $3c$; forbidding this would wrongly exclude $H$ itself.
This corrects a natural but erroneous reading of the hexagon conditions in which
\emph{all} equal arcs are forbidden; we verified numerically (\texttt{check.py})
that the present (H) exactly matches the true geodeticity test, whereas the
``all equal arcs forbidden'' variant fails already on $G'=H$.
\end{remark}

\begin{theorem}[Characterization]\label{thm:main}
For $x\in\mathbb Z_{>0}^{175}$, the homeomorph $G'(x)$ is geodetic if and only
if both \textnormal{(P)} and \textnormal{(H)} hold.
\end{theorem}

We verified Theorem~\ref{thm:main} computationally by comparing the predicate
``(P) and (H)'' against a direct BFS-based geodeticity test on more than
$10^3$ assignments (including the uniform-odd family, biased small-value
families forcing near-ties, and constructed (P)\&(H)-feasible assignments): no
discrepancy occurred (\texttt{check.py}). The proof occupies the rest of this
section. Two auxiliary facts about $H$ are used.

\begin{lemma}\label{lem:cycles}
\begin{enumerate}
\item Every pentagon of $H$ is induced (chordless). Two of its vertices are
either adjacent (split $1{+}4$) or at $H$--distance $2$ (split $2{+}3$); in the
split $2{+}3$ the $2$--arc is the unique $H$--geodesic and its middle vertex is
the common neighbour of Lemma~\ref{lem:struct}(2).
\item For a hexagon $C$ and two of its vertices $a,b$: a $1{+}5$ split means
$a,b$ adjacent; a $2{+}4$ split means $a,b$ at $H$--distance $2$ with the
$2$--arc the unique geodesic (middle vertex the common neighbour); a $3{+}3$
split means $a,b$ at $H$--distance $2$ with common neighbour \emph{off} $C$, so
that neither $3$--arc is an $H$--geodesic.
\end{enumerate}
\end{lemma}

\begin{proof}
A chord of a pentagon would create a triangle or $4$--cycle, contradicting
girth $5$; hence pentagons are induced and the cyclic distances $1,2$ are the
$H$--distances; the $2$--arc is a length-$2$ geodesic through its middle vertex.

For a hexagon: a $1{+}5$ split is an edge of $C$. A $2{+}4$ split gives a
length-$2$ $H$--path (the $2$--arc); girth $5$ forbids $a,b$ adjacent (else a
triangle with the $2$--arc), so they are at distance $2$ and the $2$--arc is the
geodesic with middle vertex the common neighbour. For a $3{+}3$ split, $a,b$
cannot be adjacent (a chord would give a $4$--cycle with a $3$--arc); by
diameter $2$ they are at distance $2$. Their unique common neighbour $w$
(Moore property) cannot lie on $C$: a $C$--vertex adjacent to both $a$ and $b$
would be at cyclic distance $1$ from each, impossible at cyclic distance $3$ on
a chordless $6$--cycle. Hence $w\notin C$, the geodesic $a\,w\,b$ has length $2$,
and the two $3$--arcs (length $3$) are not geodesics. (All three statements were
confirmed exhaustively over the $5250$ hexagons in \texttt{check.py}.)
\end{proof}

\subsection{Necessity}

Assume $G'(x)$ is geodetic. We show (P) and (H).

\emph{(H) holds.} Let $C$ be a hexagon and $a,b$ two of its branch vertices with
split $1{+}5$ or $2{+}4$, and suppose for contradiction the two arcs $A,B$ have
equal weight $w_A=w_B=:W$. By Lemma~\ref{lem:cycles}(2) the short arc (the
single edge in the $1{+}5$ case, the $2$--arc in the $2{+}4$ case) is the unique
$H$--geodesic from $a$ to $b$, hence by Lemma~\ref{lem:distbranch} the subdivided
short arc is the unique shortest $a$--$b$ path of $G'(x)$ \emph{unless} the long
arc ties it. But $A,B$ are internally disjoint $a$--$b$ paths of $G'(x)$ of equal
length $W$, and $W=w_{\text{short arc}}=d_{G'}(a,b)$ (the short arc realizes the
distance, being the geodesic, and the long arc equals it by assumption). Thus
$a,b$ have two distinct shortest paths in $G'(x)$, contradicting geodeticity.
Hence the arcs differ and (H) holds.

\emph{(P) holds.} Let $C=v_0v_1v_2v_3v_4$ be a pentagon and let
$W=\sum_{e\in C}x_e$ be its total $G'$--weight. Suppose $W$ is even. As a cycle
of $G'(x)$, $C$ has even length $W$, so there is a vertex $p$ of $C$ (internal or
branch) \emph{antipodal} to $v_0$, i.e.\ at $G'$--distance exactly $W/2$ from
$v_0$ along each of the two directions of $C$. We claim each of these two arc
routes is a shortest $v_0$--$p$ path, contradicting geodeticity.

Indeed, $p$ lies on a subdivided edge $e^*=v_iv_{i+1}$ at position $s$; by
\eqref{eq:internaldist},
$d_{G'}(v_0,p)=\min\!\big(d_{G'}(v_0,v_i)+s,\,d_{G'}(v_0,v_{i+1})+(x_{e^*}-s)\big)$.
By Lemma~\ref{lem:cycles}(1) the $H$--geodesic between any two pentagon vertices
is the corresponding short arc \emph{of $C$}, so for every pentagon vertex $v_j$
the distance $d_{G'}(v_0,v_j)$ is realized by an arc of $C$ itself (an off-$C$
competitor between two pentagon vertices closes a pentagon or hexagon with the
short arc, whose other arc is no shorter by the geodesic structure of
Lemma~\ref{lem:cycles}, applied to the already-established branch distances).
Consequently $d_{G'}(v_0,v_i)+s$ and $d_{G'}(v_0,v_{i+1})+(x_{e^*}-s)$ are the
$G'$--lengths of the two arc routes of $C$ from $v_0$ to $p$, which by the choice
of $p$ both equal $W/2$. Hence both arc routes are shortest $v_0$--$p$ paths and
$p$ is reached by two distinct shortest paths, contradicting geodeticity. So $W$
is odd and (P) holds.
\end{proof}

\subsection{Sufficiency}

Assume (P) and (H). We prove every pair of vertices of $G'(x)$ has a unique
shortest path, in two lemmas.

\begin{lemma}\label{lem:branchunique}
Under \textnormal{(P)} and \textnormal{(H)}, every pair of branch vertices
$u,v$ has a unique minimum-weight $H$--path, equal to the subdivided unique
$H$--geodesic.
\end{lemma}

\begin{proof}
By Lemma~\ref{lem:struct}, $u,v$ are at $H$--distance $1$ or $2$, with a unique
$H$--geodesic $P^\ast$ (the edge, resp.\ the $2$--path through the common
neighbour). Let $Q$ be a minimum-weight $H$--path from $u$ to $v$ with
$Q\ne P^\ast$; it suffices to show $\sum_Q x_e>\sum_{P^\ast}x_e$.

First, $|Q|\le 5$. Otherwise, since $H$ has diameter $2$ there is an $H$--path
$Q'$ from $u$ to $v$ with $|Q'|\in\{1,2\}$ and $Q'\ne Q$; if $Q'\ne P^\ast$ then
$Q'$ is a competitor of $H$--length $\le2$, hence weight at most that of the
length-$2$ path, certainly $\le \sum_Q x_e$ once $|Q|\ge6$ unless $\sum_Q$ is
unusually small --- we argue this cleanly: the minimum-weight competitor cannot
have $|Q|\ge6$, because $P^\ast$ has $|P^\ast|\le2$ and the second-shortest
$H$--path between any two vertices of $H$ has $H$--length at most $5$ (for an
edge $uv$, girth $5$ gives a length-$4$ alternative on a pentagon through $uv$,
and $7$-regularity guarantees such a pentagon exists; for a non-adjacent pair,
the geodesic $u\,w\,v$ extends to a pentagon, giving a length-$3$ alternative).
Thus there is always a competitor $Q^{\dagger}\ne P^\ast$ with $|Q^{\dagger}|\le
5$; the minimum-weight competitor $Q$ satisfies $\sum_Q\le\sum_{Q^\dagger}$, and
if $|Q|\ge6$ we may replace $Q$ by $Q^\dagger$, so we may assume $|Q|\le5$.

Now $P^\ast\cup Q$, with $Q$ of minimum weight and internally disjoint from
$P^\ast$ (any shared internal vertex would let us shorten $Q$ by splicing in part
of $P^\ast$, contradicting minimality), forms a cycle $C$ of $H$ of length
$|P^\ast|+|Q|\le 2+5=7$. By girth $5$ and the case bounds:
\begin{itemize}
\item $|P^\ast|=2,\ |Q|=3$: $C$ a pentagon ($2{+}3$ split). By (P) its odd total
weight forces $\sum_{P^\ast}\ne\sum_Q$; $P^\ast$ is the geodesic arc
(Lemma~\ref{lem:cycles}(1)), and (since at $x\equiv1$ the $2$--arc beats the
$3$--arc and arc weights are coordinate-wise monotone) $P^\ast$ is the lighter,
so $\sum_{P^\ast}<\sum_Q$.
\item $|P^\ast|=2,\ |Q|=4$: $C$ a hexagon ($2{+}4$ split). By (H) the arcs
differ; $P^\ast$ is the geodesic, hence lighter, so $\sum_{P^\ast}<\sum_Q$.
\item $|P^\ast|=1,\ |Q|=4$: $C$ a pentagon ($1{+}4$ split). By (P) the odd total
forces $x_e=\sum_{P^\ast}\ne\sum_Q$; $P^\ast$ is the edge geodesic, lighter.
\item $|P^\ast|=1,\ |Q|=5$: $C$ a hexagon ($1{+}5$ split). By (H) the arcs
differ; $P^\ast$ is the edge geodesic, lighter.
\end{itemize}
(The remaining combinatorial possibility $|P^\ast|=2,|Q|=5$ would give a
$7$--cycle; but then $Q$ is not the minimum-weight competitor, as the length-$\le5$
competitor $Q^\dagger$ above has smaller or equal weight and a shorter cycle,
reducing to a case above.) In every case $\sum_{P^\ast}x_e<\sum_Q x_e$, so
$P^\ast$ is the unique minimum-weight $H$--path. By Lemma~\ref{lem:distbranch}
its subdivision is the unique shortest $u$--$v$ path of $G'(x)$.
\end{proof}

\begin{lemma}\label{lem:pentinternal}
Under \textnormal{(P)} and \textnormal{(H)}, no internal vertex $m$ of $G'(x)$
has two shortest paths to any other vertex $t$.
\end{lemma}

\begin{proof}
Let $m$ lie on the subdivided edge $e=ab$ at distance $s$ from $a$, $0<s<x_e$.

\emph{Case $t$ a branch vertex.} By \eqref{eq:internaldist},
$d_{G'}(t,m)=\min(d_{G'}(t,a)+s,\ d_{G'}(t,b)+(x_e-s))$. By
Lemma~\ref{lem:branchunique}, $d_{G'}(t,a)$ and $d_{G'}(t,b)$ are each realized
by a unique geodesic, $P_a$ and $P_b$. A two-route tie at $m$ requires
\begin{equation}\label{eq:tie}
d_{G'}(t,a)+s=d_{G'}(t,b)+(x_e-s)=d_{G'}(t,m),
\end{equation}
with the two shortest $t$--$m$ routes being $P_a$ followed by $s$ steps into $e$,
and $P_b$ followed by $x_e-s$ steps into $e$. (If only one of the two min options
attains the minimum there is a unique route, since $P_a,P_b$ are unique; so a tie
forces both options equal, i.e.\ \eqref{eq:tie}.) Concatenating $P_a$, the edge
$e$, and the reverse of $P_b$ gives a closed walk through $t,a,b$; projecting to
$H$ and using that $a,b$ are adjacent and $P_a,P_b$ are $H$--geodesics from $t$
(of $H$--length $\le2$), the cycle $C_H:=P_a\cup\{ab\}\cup P_b$ has $H$--length
$|P_a|+1+|P_b|\le2+1+2=5$, and is a pentagon (length $5$, when
$|P_a|=|P_b|=2$ and the closure is a $5$--cycle) or a shorter cycle. By girth
$5$ a cycle of length $<5$ is impossible, so $C_H$ is a pentagon (the only way
to close with $a,b$ adjacent and $|P_a|,|P_b|\le2$ giving total $\ge5$ is
$|P_a|=|P_b|=2$, length $5$, or $|P_a|=1,|P_b|=2$ etc.; in all admissible cases
the closed walk decomposes into pentagons and hexagons of $H$). The equal-length
condition \eqref{eq:tie} makes this cycle have \emph{even} total $G'$--weight
$2\,d_{G'}(t,m)$. A pentagon of even weight is excluded by (P). The only other
admissible closure is a hexagon in which the two halves split as $1{+}5$ or
$2{+}4$ at the divergence pair $(a,b)$ (the geodesic splits), excluded by (H);
a $3{+}3$ hexagon split cannot arise here because $a,b$ are \emph{adjacent}
(split $1{+}5$) once $e=ab$ is one of the cycle edges. Hence \eqref{eq:tie} is
impossible and there is no two-route tie at $m$.

\emph{Case $t$ an internal vertex} on edge $f=cd$ at distance $r$ from $c$. By
the four-way analogue of \eqref{eq:internaldist},
$d_{G'}(t,m)$ is the minimum over the four combinations of
$\{d_{G'}(c,\cdot),d_{G'}(d,\cdot)\}$ to $\{a,b\}$ plus the offsets
$r$ or $x_f-r$ and $s$ or $x_e-s$. A two-route tie forces two of these four to
be simultaneously minimal and equal; by Lemma~\ref{lem:branchunique} each
branch-to-branch distance is uniquely realized, so the tie again produces a
closed $H$--walk through $\{c,d\}$ and $\{a,b\}$ whose two halves have equal
$G'$--weight, hence a pentagon (even weight, excluded by (P)) or a hexagon whose
geodesic-split arcs tie (excluded by (H)). If $f=e$ the only route runs along
$e$ and there is no tie. Thus no two-route tie at $m$ exists.
\end{proof}

\begin{proof}[Proof of Theorem~\ref{thm:main}, sufficiency]
By Lemma~\ref{lem:branchunique} every branch pair has a unique shortest path. By
Lemma~\ref{lem:pentinternal} no pair involving an internal vertex has two
shortest paths. Hence $G'(x)$ is geodetic.
\end{proof}

\begin{remark}[On the simplified form and the count ``$5376$'']\label{rem:simpl}
Condition (P) already implies no pentagon has two equal arcs (the two arcs sum to
an odd number, so cannot be equal); thus the pentagon non-tie constraints are
subsumed by (P). The geodetic system is therefore exactly: $1260$ pentagon
parity equations (P), together with the geodesic-split hexagon inequalities (H).
For a single hexagon the binding pairs are the $6$ edges ($1{+}5$ splits) and the
$6$ $2$--arcs ($2{+}4$ splits), i.e.\ $12$ inequalities per hexagon (the three
$3{+}3$ pairs are \emph{not} constraints). We were unable to reproduce the exact
figure ``$5376$'' from the problem statement by any natural bookkeeping
($1260$, $1260+5250=6510$, $1260+12\cdot5250$, etc.\ all differ), and we record
this discrepancy honestly: the figure in the prompt appears to be inaccurate.
The mathematically correct and computationally verified system is the one in
Definition~\ref{def:conditions}.
\end{remark}

\section{The solution set is infinite}

\begin{proposition}[Uniform odd scalings]\label{prop:infinite}
For every odd positive integer $c$, the uniform assignment $x_e=c$ for all $e$
yields a geodetic homeomorph $G'(x)$ of diameter $2c$. Hence there are infinitely
many solutions, with unbounded diameter; the solution set is not finite.
\end{proposition}

\begin{proof}
For $x_e\equiv c$, every pentagon has total weight $5c$, odd since $c$ is odd:
(P) holds. For a hexagon, the $1{+}5$ split gives arcs $c\ne5c$ and the $2{+}4$
split gives $2c\ne4c$; the $3{+}3$ split ($3c=3c$) is \emph{not} a constraint
(Remark~\ref{rem:why33excluded}). So (H) holds, and by Theorem~\ref{thm:main}
$G'(x)$ is geodetic. Its diameter is $c\cdot\mathrm{diam}(H)=2c$. As $c$ ranges
over odd positive integers we get geodetic homeomorphs of every diameter
$2,6,10,\dots$, pairwise non-isomorphic. (Verified directly for $c=1,3,5,7$ in
\texttt{check.py}.)
\end{proof}

\begin{corollary}\label{cor:illposed}
The task ``produce a complete list and exact \emph{count} of all geodetic graphs
of diameter $d\ge2$ homeomorphic to $H$'' has no finite answer: by
Proposition~\ref{prop:infinite} there are infinitely many, even up to
isomorphism. A finite enumeration is meaningful only after a finiteness
normalization, e.g.\ bounding $d\le D$, or quotienting by global rescaling
$x\mapsto\lambda x$.
\end{corollary}

\section{Algorithm}

For any prescribed bound $D$ on the diameter we enumerate exactly the
positive-integer solutions $x$ with $\max_e x_e\le D$ (these include all geodetic
homeomorphs of diameter $\le D$, since $\mathrm{diam}(G')\ge\max_e x_e$). We
exploit Theorem~\ref{thm:main} and the parity structure.

\begin{algorithm}[H]
\caption{Enumerate geodetic homeomorphs of $H$ with $\max_e x_e\le D$}
\begin{algorithmic}[1]
\Require $H$ with edge set $E$, $|E|=175$; its $1260$ pentagons $\mathcal P$ and
$5250$ hexagons $\mathcal H$; a bound $D\ge1$.
\Ensure The list $\mathcal S$ of all $x\in\{1,\dots,D\}^E$ with $G'(x)$ geodetic.
\State Build the $\mathbb F_2$ system $Ab=\mathbf 1$ in variables
$b_e=x_e\bmod2$, one row $\sum_{e\in C}b_e=1$ per pentagon.
\State Gauss-eliminate over $\mathbb F_2$: rank $r=126$, free dimension
$175-r=49$; obtain a particular solution $b^{(0)}$ and a nullspace basis.
\State $\mathcal S\gets\emptyset$.
\For{each $b$ in the affine parity space $b^{(0)}+\mathrm{span}_{\mathbb F_2}(N)$
compatible with $\{1,\dots,D\}$}
  \State $L_e\gets\{v\in\{1,\dots,D\}:v\equiv b_e\ (\mathrm{mod}\ 2)\}$.
  \State \textbf{Propagation:} for each hexagon and each $1{+}5$ or $2{+}4$ pair,
  delete domain values forced to create an arc tie; prune $b$ if a domain empties.
  \State \textbf{Backtracking search} over $\prod_e L_e$, edge order chosen to
  complete hexagons early; on each completed hexagon test its $12$
  geodesic-split inequalities (H) and backtrack on a tie.
  \State For each completed $x$, \textbf{verify} (P) on all pentagons and (H) on
  all hexagons; if both hold, add $x$ to $\mathcal S$.
\EndFor
\State \Return $\mathcal S$.
\end{algorithmic}
\end{algorithm}

\begin{proposition}[Correctness]
The algorithm returns exactly $\{x\in\{1,\dots,D\}^E:G'(x)\text{ geodetic}\}$.
\end{proposition}

\begin{proof}
By Theorem~\ref{thm:main}, $G'(x)$ is geodetic iff (P) and (H). (P) is
equivalent to $Ab=\mathbf 1$ with $b_e=x_e\bmod2$, so every solution's parity
lies in the affine space of line~4, and conversely. The search (lines 6--8)
enumerates all $x\in\prod_e L_e$ of the prescribed parity; propagation removes
only values that cannot occur in any (H)-feasible completion, so no solution is
lost; line~9 accepts $x$ iff (P) \emph{and} (H) hold. Since (P) is guaranteed by
the parity class and (H) is explicitly tested, the accepted set is exactly the
geodetic assignments with $\max_e x_e\le D$.
\end{proof}

\begin{proposition}[Termination and complexity]
The algorithm halts. It examines at most $2^{49}$ parity classes; each inner
search is over a finite product $\prod_e L_e$ with $|L_e|\le\lceil D/2\rceil$.
Each (H)-test costs $O(|\mathcal H|)=O(5250)$ and each full verification
$O(|\mathcal P|+|\mathcal H|)=O(6510)$ constant-time arc evaluations.
\end{proposition}

\begin{proof}
The parity space has dimension $49$, so line~4 iterates $\le 2^{49}$ times; each
inner search ranges over a finite set and terminates. The cost bounds follow by
counting per-iteration work.
\end{proof}

\begin{remark}[Practical feasibility --- honest assessment]\label{rem:feasible}
The pentagon parities have $\mathbb F_2$-rank $126$, leaving only $49$ free
parity bits. Each edge lies in many pentagons and hexagons, so fixing a small
spanning set of edge values propagates strongly via (P) and (H). For the smallest
normalization $D=1$ only $x\equiv1$ survives, returning $H$ itself; restricting
to constant assignments returns precisely the odd $c\le D$
(Proposition~\ref{prop:infinite}). We implemented and ran the
$O(6510)$ verifier of (P)\&(H) against a brute-force BFS geodeticity test on
$>10^3$ assignments to validate Theorem~\ref{thm:main}
(\texttt{check.py}). We do \emph{not} claim a worst-case polynomial bound for
general $D$: the naive bound $2^{49}\cdot(\lceil D/2\rceil)^{175}$ is hopeless,
and although constraint propagation collapses it dramatically in the regimes we
tested, an exact reduced-complexity bound for the hexagon-coupled search is left
open. We flag this honestly rather than overstate ``practical runtime.''
\end{remark}

\section{The enumeration result}

\begin{theorem}[Classification]\label{thm:class}
The geodetic homeomorphs $G'(x)$ of $H$ are exactly the $x\in\mathbb Z_{>0}^{175}$
satisfying
\begin{itemize}
\item[\textnormal{(P)}] $\sum_{e\in C}x_e$ odd for each of the $1260$ pentagons
$C$, and
\item[\textnormal{(H)}] for each of the $5250$ hexagons, the two arcs of every
$1{+}5$ and every $2{+}4$ split have distinct weights (the $3{+}3$ split imposes
no constraint).
\end{itemize}
The parity vectors $b=(x_e\bmod2)$ form an affine subspace of $\mathbb F_2^{175}$
of dimension $49$ (solution set of the consistent rank-$126$ system
$Ab=\mathbf1$). The all-ones assignment ($G'=H$) is a solution, as is every
uniform odd $x_e\equiv c$; hence there are infinitely many pairwise
non-isomorphic geodetic homeomorphs (distinguished by diameter $2c$). There is no
finite exact count without a finiteness normalization; under ``diameter $\le D$''
the count is finite and computed exactly by the Algorithm, while under ``constant
assignments'' the solutions are precisely $x_e\equiv c$, $c$ odd.
\end{theorem}

\begin{proof}
The characterization is Theorem~\ref{thm:main} with Remark~\ref{rem:simpl}. The
parity statement: $b\mapsto Ab$ is $\mathbb F_2$-linear of rank $126$, and
$Ab=\mathbf1$ is consistent ($b=\mathbf1$ works since each pentagon has $5\equiv1$
edges), so its solution set is affine of dimension $175-126=49$. Infinitude and
non-isomorphism are Proposition~\ref{prop:infinite}/Corollary~\ref{cor:illposed}.
\end{proof}

\begin{remark}[On $\mathrm{Aut}(H)$ and symmetry reduction]
$H$ is edge-transitive with $\mathrm{Aut}(H)\cong\mathrm P\Sigma\mathrm U(3,5)$
of order $252000$, acting on $E(H)$ and on the solution set; same-orbit solutions
give isomorphic $G'$. For a finite normalized set (e.g.\ diameter $\le D$) the
number of distinct homeomorphs is, by Burnside,
$\frac1{|\mathrm{Aut}(H)|}\sum_g|\mathrm{Fix}(g)|$, $\mathrm{Fix}(g)$ the
$D$-bounded solutions fixed by $g$. The constant-assignment family is pointwise
fixed, so each odd $c$ is one isomorphism class, consistent with distinct
diameters $2c$.
\end{remark}

\section{Summary of what is established}

\begin{enumerate}
\item \textbf{(Characterization, Theorem~\ref{thm:main}, corrected.)} $G'(x)$ is
geodetic iff every pentagon has odd weight (P) and every $1{+}5$/$2{+}4$ hexagon
split has distinct arc weights (H). The $3{+}3$ hexagon split is \emph{not} a
constraint --- correcting the natural but false ``all arcs distinct'' reading,
which already fails on $G'=H$ (Remark~\ref{rem:why33excluded}). Verified against
direct geodeticity testing in \texttt{check.py}.
\item \textbf{(Global from local, Lemmas
\ref{lem:branchunique}--\ref{lem:pentinternal}.)} Uniqueness of all geodesics
(including through internal vertices) follows from (P) and (H) via the
Moore/diameter-$2$/girth-$5$ geometry.
\item \textbf{(Infiniteness, Proposition~\ref{prop:infinite},
Corollary~\ref{cor:illposed}.)} The solution set is infinite; the literal exact
count is ill-posed and needs a finiteness normalization.
\item \textbf{(Algorithm.)} A correct, terminating enumeration for the
$D$-bounded solution set, with explicit parity structure (rank $126$, $49$ free
bits). Practical feasibility is argued but not proven in the worst case, which we
flag honestly.
\end{enumerate}

\end{document}
